\lecture{7}{2026.04.27}{}

\section{Iterative Methods}

An iterative method is to have the process of \[
    x_1, x_2, \ldots, \to x^*
\]
We hope to find \[
    Ax^* = b
\]
This type of method should be faster than Gaussian elimination $O(n^3)$. If $x_k \to x_{k+1}$ takes $O(n^r)$, then we can have $O(ln^r)$, where $l$ is the number of iterations. There are two properties should be satisfied for an iterative method:
\begin{itemize}
    \item Numerical stability
    \item Sparsity (Fill-in)
\end{itemize}

\subsection{Jacobi Method}

\begin{eg}
    Consider the system of equations $Ax = b$
\end{eg}

We have 
\begin{eqnarray*}
x_1 &=& (b_1 - a_{12}x_2 - a_{13}x_3)/a_{11}\\
x_2 &=& (b_2 - a_{21}x_1 - a_{23}x_3)/a_{22}\\
x_3 &=& (b_3 - a_{31}x_1 - a_{32}x_2)/a_{33}
\end{eqnarray*}
The Jacobi method is to update $x_k$ to $x_{k+1}$ by 
\begin{eqnarray*}
(x_{k+1})_1 &=& (b_1 - a_{12}(x_k)_2 - a_{13}(x_k)_3)/a_{11}\\
(x_{k+1})_2 &=& (b_2 - a_{21}(x_k)_1 - a_{23}(x_k)_3)/a_{22}\\
(x_{k+1})_3 &=& (b_3 - a_{31}(x_k)_1 - a_{32}(x_k)_2)/a_{33}
\end{eqnarray*}

In the general case, we should do the procedure of 

\begin{codebox}
\Procname{$\proc{Jacobi}(A, b, x^{(0)})$}
\li \For $i = 1 \twodots n$
\Do
\li $(x_{k+1})_i =  ( b_i - \sum_{j=1}^{i-1} a_{ij} (x_k)_j - \sum_{j=i+1}^n a_{ij} (x_k)_j )/ a_{ii}$
\End
\end{codebox}

\begin{definition}[Jacobi Method]
    The Jacobi method is to update $x_k$ to $x_{k+1}$ by 
    \[
       (x_{k+1})_i = \frac{1}{a_{ii}} \left( b_i - \sum_{j=1}^{i-1} a_{ij} (x_k)_j - \sum_{j=i+1}^n a_{ij} (x_k)_j \right)
    \]
\end{definition}

\subsection{Gauss-Seidel Method}

Now consider Gauss-Seidel method, which is to update $x_k$ to $x_{k+1}$ using the recently updated values. We have
\begin{eqnarray*}
(x_{k+1})_1 &=& (b_1 - a_{12}(x_k)_2 - a_{13}(x_k)_3)/a_{11}\\
(x_{k+1})_2 &=& (b_2 - a_{21}(x_{k+1})_1 - a_{23}(x_k)_3)/a_{22}\\
(x_{k+1})_3 &=& (b_3 - a_{31}(x_{k+1})_1 - a_{32}(x_{k+1})_2)/a_{33}
\end{eqnarray*}

For the general case, we should do the procedure of
\begin{codebox}
\Procname{$\proc{Gauss-Seidel}(A, b, x^{(0)})$}
\li \For $i = 1 \twodots n$
\Do
\li $(x_{k+1})_i =  ( b_i - \sum_{j=1}^{i-1} a_{ij} (x_{k+1})_j - \sum_{j=i+1}^n a_{ij} (x_k)_j )/ a_{ii}$
\End
\end{codebox}

\begin{definition}
    The Gauss-Seidel method is to update $x_k$ to $x_{k+1}$ by 
    \[
        (x_{k+1})_i = \frac{1}{a_{ii}} \left( b_i - \sum_{j=1}^{i-1} a_{ij} (x_{k+1})_j - \sum_{j=i+1}^n a_{ij} (x_k)_j \right)
    \]
\end{definition}

The iterative method can be divergent, 

\begin{eg}
    \begin{equation*}
    A =
    \begin{pmatrix}
      1 & 2 & 2 \\ 
    2 & 1 & 2 \\
    2 & 2 & 1
    \end{pmatrix},
    b =
    \begin{pmatrix}
      5 \\ 5 \\ 5
    \end{pmatrix},
    \mbox{sol} =
    \begin{pmatrix}
      1 \\ 1 \\ 1
    \end{pmatrix}
    \end{equation*}
\end{eg}

The determinant of $A$ is $-9$, so $A$ is invertible. However, the iterative method will diverge 
\begin{equation*}
x_0 = 
\begin{pmatrix}
  0 \\ 0 \\ 0
\end{pmatrix},\
x_1 =
\begin{pmatrix}
  5\\ 5 \\ 5
\end{pmatrix},\
x_2 
= 
\begin{pmatrix}
5 - 10 - 10  \\
5 - 10 - 10  \\
5 - 10 - 10  
\end{pmatrix}
= 
\begin{pmatrix}
  -15 \\ -15 \\ -15
\end{pmatrix},\
    x_3 
=
\begin{pmatrix}
  5 + 30 + 30 \\
  5 + 30 + 30 \\
  5 + 30 + 30 
\end{pmatrix}
= 
\begin{pmatrix}
  65 \\ 65\\ 65
\end{pmatrix}
\end{equation*}

Here is the convergence analysis of the iterative method. The goal is to find when the iterative method will get the solution of $Ax = b$.
\begin{eqnarray*}
a_{11} x_1 &=& 
b_1 - a_{12}x_2 - a_{13}x_3\\
a_{22} x_2 &=& b_2 - a_{21}x_1 - a_{23}x_3\\
a_{33} x_3 &=& b_3 - a_{31}x_1 - a_{32}x_2
\end{eqnarray*}
Hence, we can rewrite the above equations as \begin{equation*}
  \begin{pmatrix}
    a_{11} && \\
& a_{22} & \\
&& a_{33}
  \end{pmatrix}
  \begin{pmatrix}
    x_1 \\ x_2 \\ x_3
  \end{pmatrix}
=
\begin{pmatrix}
  0 & -a_{12} & -a_{13} \\
 -a_{21} & 0 & -a_{23} \\
-a_{31} & -a_{32} & 0
\end{pmatrix}
  \begin{pmatrix}
    x_1 \\ x_2 \\ x_3
  \end{pmatrix}
 + b
\end{equation*}

If \[
    M = \begin{pmatrix}
    a_{11} && \\
& a_{22} & \\
&& a_{33}
  \end{pmatrix}, N = \begin{pmatrix}
  0 & -a_{12} & -a_{13} \\
 -a_{21} & 0 & -a_{23} \\
-a_{31} & -a_{32} & 0
\end{pmatrix}
\]
then \[
    A = M - N, \quad Mx_{k+1} = Nx_{k} + b
\]

Here we define the spectral radius

\begin{definition}[Spectral Radius]
    The spectral radius of a matrix $A$ is \[
    \rho(A) = \max\{|\lambda|: \lambda \in \lambda(A)\}
    \]
    where $\lambda(A)$ is the set of eigenvalues of $A$.
\end{definition}

Then we have the following theorem 
\begin{theorem}
    Assume \[
        A = M - N
    \]
    where \[
        A, M \mbox{are non-singular and } \rho(M^{-1}N) < 1
    \]
    Then \[
        M x_{k+1} = Nx_k + b
    \]
    leads to convergence of $\{x_k\}$ to $A^{-1}b$ for any initial guess $x_1$.
\end{theorem}
\begin{proof}
    We have \[
        Ax = b \implies Mx = Nx + b
    \]
    Get the iterative method in \[
        M(x_{k+1} - x) = N(x_k - x)
    \]
    Hence, \[
        x_{k+1} - x = M^{-1}N(x_k - x)
    \]
    Get the general formula 
    \[
        x_{k+1} - x = (M^{-1}N)^{k}(x_1 - x)
    \]
    With the condition $\rho(M^{-1}N) < 1$, we have \[
        \lim_{k \to \infty} (M^{-1}N)^k = 0
    \]

\begin{note}
    The reason of condition \[
        \rho(M^{-1}N) < 1 \implies \lim_{k \to \infty} (M^{-1}N)^k = 0
    \] 
    is quite complicated, we omit the proof here.
\end{note}

    Hence, \[
        \lim_{k \to \infty} x_{k+1} - x = 0
    \]
\end{proof}

\subsection{Theoretical Explanation of Gauss-Seidel Method}

Now we try to give Gauss-Seidel method a theoretical explanation. We construct the following opitmization problem \[
    \min_{x} f(x) = \frac{1}{2} x^T A x - b^T x
\]
which is the same as solving $Ax = b$ if $A$ is a symmetric positive definite. We let \[
    f(x) = \frac{1}{2} x^T A x - b^T x
\]
then \[
    \nabla f(x) \equiv \begin{pmatrix}
        \frac{\partial f}{\partial x_1} \\
        \vdots \\
        \frac{\partial f}{\partial x_n}
    \end{pmatrix} = Ax - b
\]

The derivation will be \begin{align*}
    x^T A x &= \sum_{i=1}^n x_i (Ax)_i \\[3pt]
    &= \sum_{i=1}^n \sum_{j=1}^n x_i a_{ij} x_j \\[3pt]
    &= \underbrace{x_1  A_{11} x_1 + \cdots + x_1 A_{1n} x_n}_{\text{first row}} + \underbrace{x_2 A_{21} x_1 + \cdots + x_2 A_{2n} x_n}_{\text{second row}} + \cdots + \underbrace{x_n A_{n1} x_1 + \cdots + x_n A_{nn} x_n}_{n\text{-th row}}
\end{align*}
Therefore, \begin{align*}
    \frac{\partial x^T A x}{\partial x_1} &= (2 A_{11} x_1 + A_{12} x_2 + \cdots + A_{1n} x_n) + x_2 A_{21} + \cdots + x_n A_{n1} \\[3pt]
    &= 2 (A_{11} x_1 + A_{12} x_2 + \cdots + A_{1n} x_n) \tag{since $A$ is symmetric}
\end{align*}

\begin{note}
    Coordinate-wise minimization is a classic way for solving optimization problems. Sequentially optimizing one variable at a time \begin{gather*}
        \min_{x_1} f(x_1, x^k_2, \ldots, x^k_n)\\
        \min_{x_2} f(x^{k+1}_1, x_2, \ldots, x^k_n)\\
        \min_{x_3} f(x^{k+1}_1, x^{k+1}_2, x_3, \ldots, x^k_n)\\
        \vdots
  \end{gather*}
\end{note}

Assume $x_i$ is updated \[
    \min_{d} \frac{1}{2} (x + d e_i)^T A (x + d e_i) - b^T (x + d e_i)
\]

Next, \[
    \min_{d} \frac{1}{2} d^2 A_{ii} + d (Ax)_i - d b_i
\]
\[
    A_{ii} d + (Ax)_i - b_i = 0
\]
The minimum is achieved at \[
    d = \frac{b_i - (Ax)_i}{A_{ii}}
\]
So, \[
    x_i + d \leftarrow \frac{b_i - \sum_{j:j \neq i} A_{ij} x_j}{A_{ii}}
\]
We can get the Gauss-Seidel update rule.

\section{Conjugate Gradient Method}

This method can use only symmetric positive definite matrices. We first disscuss the \textbf{steepest descent method}. 

\subsection{Steepest Descent Method}

We construct the following optimization problem \[
    \min_{x} f(x) = \frac{1}{2} x^T A x - b^T x
\]
Using taylor expansion, we have \[
    f(x) = f(x_k) + \nabla f(x_k)^T (x - x_k) + \frac{1}{2} (x - x_k)^T \nabla^2 f(x_k) (x - x_k) + \cdots
\]
In this method, we omit the higher order terms and only consider the first two terms. Hence, we have \[
    f(x) \approx f(x_k) + \nabla f(x_k)^T (x - x_k)
\]
The first term $f(x_k)$ is a constant, so we only need to minimize $\nabla f(x_k)^T (x - x_k)$. If \[
    \min_{\|x - x_k\| = 1} \nabla f(x_k)^T (x - x_k)
\]
then \[
    x - x_k = - \frac{\nabla f(x_k)}{\|\nabla f(x_k)\|}
\]
is the direction of steepest descent, which is the unit vector of $-\nabla f(x_k)$.
\begin{note}
    We need to \[
        \| x - x_k \| = 1
    \]
    otherwise, we can make $x$ go to infinity and make $\nabla f(x_k)^T (x - x_k)$ go to negative infinity.
\end{note}

Now we have \[
    \nabla f(x_k) = Ax_k - b
\]
Let \[
    r = - \nabla f(x_k) = b - Ax_k, \quad x = x_k
\]
We minimize along the direction of $r$ 
\begin{eqnarray*}
&& \min_{\alpha}
\frac{1}{2} (x + \alpha r)^T A(x+ \alpha r)
 - (x + \alpha r)^T b \\  
&& \min_{\alpha}
\frac{1}{2} \alpha^2 r^T Ar
 + \alpha r^T A x - 
\alpha r^T b \\  
&& \min_{\alpha}
\frac{1}{2} \alpha^2 r^T Ar
 + \alpha r^T( A x - b) 
\end{eqnarray*}

Then, it becomes an one-dimensional optimization problem
\begin{eqnarray*}
&& \alpha r^T Ar + r^T(Ax - b) = 0 \\[6pt]
&& \alpha 
= \frac{r^T(b - Ax)}{r^T A r}
\end{eqnarray*}

\begin{note}
    $r^T A r \neq 0$ since $A$ is positive definite and $r \neq 0$.
\end{note}

Now $r = b - Ax$, we have \[
    \alpha = \frac{(b - Ax)^T (b - Ax)}{(b - Ax)^T A (b - Ax)} = \frac{r^T r}{r^T A r}
\]

Hence, we can get the update rule of the steepest descent method 
\begin{codebox}
\Procname{$\proc{SteepestDescent}(A, b)$}
\li $k = 0,\ x_0 = 0,\ r_0 = b$
\li \While $r_k \neq 0$
\Do
\li $k = k + 1$
\li $\alpha_k = \frac{r^T_{k-1} r_{k-1}}{r^T_{k-1} A r_{k-1}}$
\li $x_k \leftarrow x_{k-1} + \alpha_k r_{k-1}$
\li $r_k \leftarrow b - A x_k$
\End
\end{codebox}

However, the steepest descent method may \red{converge very slowly}, becasue $r$ is not a good direction to go. So we have to construct a better direction. Before that, we first construct the general search direction method.

\begin{definition}[Search Direction]
    Suppose we have $x_k$ obtained by \[
        \min_{\alpha} f(x_{k-1} + \alpha p_k)
    \]
    where $p_k$ is the search direction.
\end{definition}

Then, \begin{align*}
    f(x_{k-1} + \alpha p_k) &= \frac{1}{2} (x_{k-1} + \alpha p_k)^T A (x_{k-1} + \alpha p_k) - b^T (x_{k-1} + \alpha p_k)\\
    &= C + \alpha p_k^T (A x_{k-1} - b) + \frac{1}{2} \alpha^2 p_k^T A p_k
\end{align*}
So, \[
    \alpha = - \frac{p_k^T (r_{k-1})}{p_k^T A p_k}
\]
Hence, we have the following gereral search direction method
\begin{codebox}
\Procname{$\proc{SearchDirection}(A, b)$}
\li $k = 0,\ x_0 = 0,\ r_0 = b$
\li \While $r_k \neq 0$
\Do
\li $k = k + 1$
\li \red{Choose a search direction $p_k$ such that $p_k^T r_{k-1} \neq 0$}
\li $\alpha_k = - \frac{p_k^T r_{k-1}}{p_k^T A p_k}$
\li $x_k \leftarrow x_{k-1} + \alpha_k p_k$
\li $r_k \leftarrow b - A x_k$
\End
\end{codebox}

In this setting, \[
    x_k = x_{0} + \mbox{span}\{p_1, \ldots, p_k \}
\]
The question become how to choose the search direction $p_k$?

\newpage

\subsection{Conjugate Gradient Method}

For selecting the search direction, we hope that $p_1, \ldots, p_k$ are linearly independent and \begin{equation}
    x_k = \arg \min_{x \in x_0 + \text{span}\{p_1, \ldots, p_k\}} f(x) \label{eq:arg}
\end{equation}

With the condition \ref{eq:arg}, \[
    \mbox{span}\{p_1, \ldots, p_n\} = \mathbb{R}^n
\]
so \[
    x_n = \arg \min_{x \in \mathbb{R}^n} f(x)
\]
Then \[
    Ax_n = b
\]
and the procedure will converge in at most $n$ iterations. To maintain the property \ref{eq:arg}, we need to let \[
    x_k = x_0 + P_{k-1} y + \alpha p_k
\]
where \[
    P_{k-1} = (p_1, \ldots, p_{k-1}), \quad y \in \mathbb{R}^{k-1}, \quad \alpha \in \mathbb{R}
\]
From the optimality condition, we have \begin{align*}
    f(x_k) &= \frac{1}{2} (x_0 + P_{k-1} y + \alpha p_k)^T A (x_0 + P_{k-1} y + \alpha p_k) - b^T (x_0 + P_{k-1} y + \alpha p_k) \\
    &= \frac{1}{2} (x_0 + P_{k-1} y)^T A (x_0 + P_{k-1} y) - b^T (x_0 + P_{k-1} y) + \alpha p_k^T A (x_0 + P_{k-1} y) - b^T (\alpha p_k)+ \frac{\alpha^2}{2}  p_k^T A p_k \\ 
    &= f(x_0 + P_{k-1} y) + \alpha p_k^T A P_{k-1} y - \alpha p_k^T r_0 + \frac{\alpha^2}{2} p_k^T A p_k
\end{align*}

In this way, sloving \[
    \min_{x \in x_0 + \text{span}\{p_1, \ldots, p_k\}} f(x) = \min_{y, \alpha} f(x_0 + P_{k-1} y + \alpha p_k)
\]
is difficult because the term \[
    \alpha p_k^T A P_{k-1} y
\]
involves both $y$ and $\alpha$.
To make the problem easier, if \[
    p_k \in \mbox{span}\{Ap_1, \ldots, Ap_{k-1}\}^\perp
\]
then \[
    p_k^T A P_{k-1} y = 0
\]
and \[
    \min_{x \in x_0 + \text{span}\{p_1, \ldots, p_k\}} f(x) = \min_{y} f(x_0 + P_{k-1} y) + \min_{\alpha} 
    (- \alpha p_k^T r_0 + \frac{\alpha^2}{2} p_k^T A p_k)
\]
Then we have two independent optimization problems. Therefore, we require \red{$p_k$ to be $A$-conjugate to $p_1, \ldots, p_{k-1}$}, i.e. \[
    p_k^T A p_j = 0, \quad \forall j < k
\]
By induction, we already have \[
    x_{k-1} = \arg \min_{y} f(x_0 + P_{k-1} y)
\]
The solution of the second optimization problem is \[
    \alpha = \frac{p_k^T r_0}{p_k^T A p_k}
\]
\newpage
Because of the $A$-conjugacy, we have \begin{align*}
    p_k^T r_{k-1} &= p_k^T (b - A x_{k-1}) \\
    &= p_k^T (b - A(x_0 + P_{k-1} y_{k-1})) = p_k^T r_0
\end{align*}
We have \[
    \alpha_k = \frac{p_k^T r_0}{p_k^T A p_k} = \frac{p_k^T r_{k-1}}{p_k^T A p_k}
\]
and \[
    x_k = x_{k-1} + \alpha_k p_k
\]
We can get the following algorithm for the conjugate gradient method
\begin{codebox}
\Procname{$\proc{ConjugateGradient}(A, b)$}
\li $k = 0,\ x_0 = 0,\ r_0 = b$
\li \While $r_k \neq 0$
\Do
\li $k = k + 1$
\li \red{Choose any $p_k \in \mbox{span}\{Ap_1, \ldots, Ap_{k-1}\}^\perp$ such that $p_k^T r_{k-1} \neq 0$}
\li $\alpha_k = \frac{p_k^T r_{k-1}}{p_k^T A p_k}$
\li $x_k \leftarrow x_{k-1} + \alpha_k p_k$
\li $r_k \leftarrow b - A x_k$
\End
\end{codebox}

We still not pick $p_k$ yet, one common way is to minimize the distance between $p_k$ and $r_{k-1}$.

\begin{note}
    Because $r_{k-1}$ is the direction of steepest descent \[
        \nabla f(x_{k-1}) = A x_{k-1} - b = - r_{k-1}
    \]
\end{note}

The algorithm is become 
\begin{codebox}
\Procname{\proc{Conjugate-Gradient}(A, b)}
\li $k \gets 0,\ x_0 \gets 0,\ r_0 \gets b$
\li \While $r_k \neq 0$ \Do
\li     $k \gets k + 1$
\li     \If $k = 1$ \Then
\li         $p_1 \gets r_0$
\li     \Else
\li         \red{Let $p_k$ minimize $\|p_k - r_{k-1}\|_2$ overall vectors in $p \in \mbox{span}\{Ap_1, \ldots, Ap_{k-1}\}^\perp$}
\End
\li     $\alpha_k \gets \frac{p_k^T r_{k-1}}{p_k^T A p_k}$
\li     $x_k \gets x_{k-1} + \alpha_k p_k$
\li     $r_k \gets b - A x_k$
\End
\end{codebox}

We have the following lemma for choosing $p_k$.
\begin{lemma} \label{lemma:pk}
    If $p_k$ minimizes $\|p_k - r_{k-1}\|_2$ overall vectors in $p \in \mbox{span}\{Ap_1, \ldots, Ap_{k-1}\}^\perp$, then \begin{equation}
        p_k = r_{k-1} - A P_{k-1} z_{k-1} \label{eq:pk}
    \end{equation}
    where $z_{k-1}$ solves \begin{equation}
        \min_{z} \|r_{k-1} - A P_{k-1} z\|_2 \label{eq:zk}
    \end{equation}
    $P_{k} = (p_1, \ldots, p_{k})$: an $n \times k$ matrix.
\end{lemma}
\begin{proof}
    Let $z_{k-1}$ be the solution of (\ref{eq:zk}), then and define \[
        p = r_{k-1} - A P_{k-1} z_{k-1}
    \]
    The optimality condition of (\ref{eq:zk}) is \[
        P_{k-1}^T A^T (r_{k-1} - A P_{k-1} z_{k-1}) = 0
    \]
    Thus \[
        p^T A^T p_i = 0, \quad \forall i < k
    \]
    Moreover, $p$ is the orthogonal projection of $r_{k-1}$ onto $S = \mbox{span}\{Ap_1, \ldots, Ap_{k-1}\}^\perp$, so $p = p_k$ is the solution of \[
        \min_{p \in S} \|p - r_{k-1}\|_2
    \]
\end{proof}

\begin{note}
    $\nabla f (x) = 0$ is the optimality condition of the optimization problem \[
        f(z) = \|r_{k-1} - A P_{k-1} z\|_2^2
    \]
    $$\nabla_z f(z) = -2 (A P_{k-1})^T r_{k-1} + 2 (A P_{k-1})^T (A P_{k-1}) z = 0$$
\end{note}

\begin{figure}[H]
    \centering
    \begin{tikzpicture}[
    x={(-0.8cm,-0.4cm)}, y={(1cm,0cm)}, z={(0cm,1cm)},
    >=Stealth,
    vector/.style={->, thick},
    plane/.style={fill=blue!5, draw=gray, opacity=0.7}
]
    \coordinate (O) at (2,1,0);
    \coordinate (Proj) at (2, 3, 0);
    \coordinate (R_tip) at (2, 3, 2.5);

    \draw[plane] (-1,-1,0) -- (4,0,0) -- (5,4,0) -- (0, 3, 0) -- cycle;
    \node[anchor=north] at (2,-2,1) {$S = \text{span}\{Ap_1, \dots, Ap_{k-1}\}$};

    \draw[vector, blue] (O) -- (Proj) node[midway, below] {$AP_{k-1}z$};
    
    \draw[vector, black] (O) -- (R_tip) node[midway, above left] {$r_{k-1}$};

    \draw[vector, red] (Proj) -- (R_tip) node[midway, right] {$p_k \perp S$};

    \begin{scope}[canvas is yz plane at x=2]
        \draw[thin] (3, 0.3) -- (2.7, 0.3) -- (2.7, 0);
    \end{scope}
    \begin{scope}[canvas is xz plane at y=3]
        \draw[thin] (2, 0.3) -- (1.7, 0.3) -- (1.7, 0);
    \end{scope}

\end{tikzpicture}
    \caption{The geometric interpretation of the conjugate gradient method.}
\end{figure}


\subsection{Properties of Conjugate Gradient Method}

\begin{theorem}
    After $j$ iterations, we have 
    \begin{eqnarray}
        r_j & = & r_{j-1} - \alpha_j Ap_j \nonumber \\
        P_j^T r_j &= & 0 \label{eq:Pr}\\
        \mbox{span}\{p_1, \ldots, p_j\}
        & = & 
        \mbox{span}\{r_0, \ldots, r_{j-1}\} \nonumber \\
        & = & \mbox{span}\{b, Ab, \ldots,
        A^{j-1}b\} \nonumber
    \end{eqnarray}
    \[
        r_i^T r_j = 0 \mbox{ for all } i \neq j
    \]
\end{theorem}
\begin{proof}
    We only proof the first property, the others can be proved by induction. \begin{align*}
        r_j &= b - A x_j \\
        &= b - A(x_{j-1} + \alpha_{j-1} p_{j-1}) \\
        &= r_{j-1} - \alpha_j A p_j
    \end{align*}
\end{proof} 

From this theorem, we can see that $r_i$, $r_j$ are orthogonal for $i \neq j$.

\begin{recall}
    lemma \ref{lemma:pk}
    \[
        p_k = r_{k-1} - A P_{k-1} z_{k-1}
    \]
\end{recall}

\red{Now we want to find $z_{k-1}$} because $A, P_{k-1}$ are known. We have $z_{k-1}$ is a vector with length $k-1$, \[
    z_{k-1} = \begin{pmatrix}
        \mathbf{w} \\ \mu
    \end{pmatrix}, \quad \mathbf{w} \in \mathbb{R}^{k-2}, \mu \in \mathbb{R}
\]
\begin{eqnarray}
    p_k & = & r_{k-1} - AP_{k-1} z_{k-1} \label{eq:rApz}\\
& = &r_{k-1} - AP_{k-2}\mathbf{w} - \mu Ap_{k-1} \nonumber \\
& = & (1 + \frac{\mu}{\alpha_{k-1}}) r_{k-1} + s_{k-1}
\label{eq:pk}
\end{eqnarray}

From the thorem, we have \[
    r_{k-1} = r_{k-2} - \alpha_{k-1} A p_{k-1}
\]

\begin{definition}
    We call $s_{k-1}$ the \textbf{correction term} of $p_k$. \begin{equation}
        \begin{split}
        s_{k-1}
        & =
        -\frac{\mu}{\alpha_{k-1}} r_{k-1}
        - AP_{k-2}\mathbf{w}
        - \mu A p_{k-1}\\
        &
        = -\frac{\mu}{\alpha_{k-1}} r_{k-2}
        - AP_{k-2}\mathbf{w}  
        \end{split}
        \label{eq:sk-1}
    \end{equation}
\end{definition}

Then, we have \[
    r_i^T r_{j} = 0, \quad \forall i \neq j
\]
and from the theorem, we have \begin{equation*}
  \begin{split}
    AP_{k-2}\mathbf{w} & \in
    \text{span}\{Ap_1, \ldots,  A p_{k-2}\} \\
    & =     \text{span}\{Ab, \ldots,  A^{k-2}b\}\\
& \subset
\text{span}\{r_0, \ldots, r_{k-2}\}
\end{split}
\end{equation*}

Hence, $r_{k-1}^T (AP_{k-2}\mathbf{w}) = 0$
\begin{equation}
    s_{k-1}^T r_{k-1} = 0 \label{eq:rw}
\end{equation}

\begin{recall}
    From the earlier results, our goal is to find $z_{k-1}$ such that \[
        \| r_{k-1} - A P_{k-1} z\|_2
    \]
    is minimized.
\end{recall}

The reason of minimizing $\| r_{k-1} - A P_{k-1} z\|_2$ instead of \begin{equation}
    \| p - r_{k-1} \|_2, \quad p \in \mbox{span}\{Ap_1, \ldots, Ap_{k-1}\}^\perp \label{eq:not}
\end{equation}
is that the optimization problem (\ref{eq:not}) is constrained.

From (\ref{eq:rApz}) and (\ref{eq:pk}), we want to select $\mu$ an $\mathbf{w}\ (s)$ to minimize \[
    \left\| \left( 1 + \frac{\mu}{\alpha_{k-1}} \right) r_{k-1} + s_{k-1} \right\|^2
\]
From (\ref{eq:rw}), we have the equivalent \[
    \left\| \left( 1 + \frac{\mu}{\alpha_{k-1}} \right) r_{k-1} \right\|^2 + \left\| - \frac{\mu}{\alpha_{k-1}} r_{k-2} - AP_{k-2}\mathbf{w} \right\|^2
\]
If an optimal solution is $(\mu^*, \mathbf{w}^*)$, then \[
    \left\| - \frac{\mu^*}{\alpha_{k-1}} r_{k-2} - AP_{k-2}\mathbf{w}^* \right\|^2 = \left| -\frac{\mu^*}{\alpha_{k-1}} \right|^2 \left\| r_{k-2} - AP_{k-2} \boxed{\frac{\mathbf{w}^*}{-\mu^* / \alpha_{k-1}}} \right\|^2
\]
(The boxed term is the optimization problem to find $p$ at $k-1$ steps) and \[
    -\mu^* / \alpha_{k-1}
\]
must be the optimal solution of \[
    \min_{z} \|r_{k-2} - A P_{k-2} z\|_2
\]

From the earlier lemma, the solution of \[
    \min_{z} \|r_{k-2} - A P_{k-2} z\|_2
\]
is \[
    p_{k-1} = r_{k-2} - A P_{k-2} z_{k-2}
\]
Therefore, $s_{k-1}$ is parallel to $p_{k-1}$, from (\ref{eq:pk}), we have \[
    p_k \in \mbox{span}\{r_{k-1}, p_{k-1}\}
\]

Now, \red{assume} \begin{equation}
    p_k = r_{k-1} + \beta_{k} p_{k-1} \label{eq:pk-beta}
\end{equation}
This assumption is reasonable because we will adjust the step size $\alpha_k$ to make the direction of $p_k$ correct. Therefore, finding a \red{parallel} direction to the real solution of $\min \|p-r_{k-1}\|_2$ is enough. \\

From the $A$-conjugacy of $p_i$ and \ref{eq:Pr}, we respectively have \[
    p_{k-1}^T A p_{k} = 0,\quad p_{k-1}^T r_{k-1} = 0
\]
With (\ref{eq:pk-beta}) we have \[
    Ap_k = A r_{k-1} + \beta_{k} A p_{k-1}
\]
\[
    0 = p_{k-1}^T A p_k = p_{k-1}^T A r_{k-1} + \beta_{k} p_{k-1}^T A p_{k-1}
\]
So, \[
    \beta_{k} = - \frac{p_{k-1}^T A r_{k-1}}{p_{k-1}^T A p_{k-1}}
\]
which is a closed-form solution for $\beta_{k}$.
\[
    \alpha_k = \frac{p_k^T r_{k-1}}{p_k^T A p_k} = \frac{(r_{k-1} + \beta_k p_{k-1})^T r_{k-1}}{p_k^T A p_k} = \frac{r_{k-1}^T r_{k-1}}{p_k^T A p_k}
\]
The algorithm is become
\begin{codebox}
\Procname{\proc{Conjugate-Gradient}(A, b)}
\li $k \gets 0,\ x_0 \gets 0,\ r_0 \gets b$
\li \While $r_k \neq 0$ \Do
\li     $k \gets k + 1$
\li     \If $k = 1$ \Then
\li         $p_1 \gets r_0$
\li     \Else
\li         \red{$\beta_k \gets - \frac{p_{k-1}^T A r_{k-1}}{p_{k-1}^T A p_{k-1}}$}
\li         \red{$p_k \gets r_{k-1} + \beta_k p_{k-1}$}
\End
\li     \red{$\alpha_k \gets \frac{r_{k-1}^T r_{k-1}}{p_k^T A p_k}$}
\li     $x_k \gets x_{k-1} + \alpha_k p_k$
\li     $r_k \gets b - A x_k$
\End
\end{codebox}

However, the above algorithm is not efficient because we need to compute three matrix-vector products in each iteration \[
    A p_{k-1}, A x_k, A r_{k-1}
\]